\chapter{Análisis estructural de la red}

    En el capítulo anterior partimos de los ficheros \textit{GTFS} y, mediante un proceso de limpieza, validación y consolidación, construimos un grafo que representa la red de transporte público. En él, cada nodo agrupa un conjunto de paradas cercanas, mientras que las aristas reflejan los trayectos efectivos entre ellas según las rutas registradas.

    Una vez obtenido este grafo, el paso siguiente es examinar su estructura para comprender cómo se organiza y qué tan eficiente o robusta es. Existen múltiples enfoques posibles; en este capítulo se adoptan cuatro: el cálculo de métricas globales del grafo, el análisis de centralidades a nivel de nodo, la detección de comunidades y la evaluación de robustez ante fallas o ataques dirigidos.

    \section{Métricas globales}
    \section{Centralidades: algoritmos y parámetros}
    \section{Detección de comunidades}
    \section{Análisis de robustez: escenarios y procedimiento}
    \section{Parámetros, software y reproducibilidad}
